% operator-stability-v1 — publication template (writeup-aligned)
% Slots filled by scripts/generate_operator_paper.py
\documentclass[12pt]{article}
\usepackage[letterpaper,top=2cm,bottom=2cm,left=3cm,right=3cm]{geometry}
\usepackage[T1]{fontenc}
\usepackage[utf8]{inputenc}
\usepackage{lmodern}
\usepackage{microtype}
\usepackage{amsmath,amssymb,amsthm,mathtools}
\usepackage{booktabs}
\usepackage{xcolor}
\usepackage{enumitem}
\usepackage{xurl}
\usepackage[colorlinks=true,linkcolor=blue,urlcolor=blue,citecolor=blue]{hyperref}

\Urlmuskip=0mu plus 1mu
\urlstyle{tt}
\emergencystretch=2em

\setlength{\parskip}{0.4em}
\setlength{\parindent}{1.2em}
\setlist[itemize]{leftmargin=1.4em,itemsep=0.25em,topsep=0.35em}
\setlist[enumerate]{leftmargin=1.4em,itemsep=0.25em,topsep=0.35em}

\theoremstyle{plain}
\newtheorem{theorem}{Theorem}
\newtheorem{lemma}[theorem]{Lemma}
\theoremstyle{definition}
\newtheorem{definition}[theorem]{Definition}
\newtheorem{example}[theorem]{Example}
\theoremstyle{remark}
\newtheorem{remark}[theorem]{Remark}

% Breakable monospace for Lean paths / identifiers (xurl \path)
\newcommand{\leanpath}[1]{%
  {\raggedright\urlstyle{tt}\path{#1}}}


\definecolor{fundbox}{RGB}{245,245,248}
\definecolor{fundrule}{RGB}{180,180,190}

\title{Weighted tournament winner Margin Preservation under Bounded Score Perturbations}
\author{Research Architecture Runtime\\[0.25em]
{\normalsize\texttt{operator/weighted-tournament-winner}}}
\date{2026-07-29}

\begin{document}
\maketitle

\begin{abstract}
\leanpath{weighted-tournament-winner} selects a unique maximizer. Margin $\gamma(s)>2\varepsilon$ is necessary and sufficient for unique-max invariance under $\|\delta\|_\infty\le\varepsilon$. Lean \texttt{LEAN\_FULL} (reduction to \texttt{Argmax.Margin}).
\end{abstract}

\noindent
\begin{center}
\fcolorbox{fundrule}{fundbox}{%
  \begin{minipage}{0.92\linewidth}
  \centering
  \textbf{This theorem is:} \texttt{Reduction}
  \end{minipage}}
\end{center}
\vspace{0.6em}

\section{Problem}
{\raggedright\sloppy
The \leanpath{weighted-tournament-winner} operator returns the unique maximizer index $i^\star$ of scores $s\in\mathbb{R}^m$ ($m\ge 2$), up to the operator's definitional presentation (heap top, tournament winner, masked max, etc.).
\par}

\section{Stability notion}
Perturbations: $\|\delta\|_\infty\le\varepsilon$. Let $\gamma(s)=s_{i^\star}-\max_{j\neq i^\star}s_j$. Stability: $\gamma(s)>2\varepsilon$ implies $i^\star$ remains the unique maximizer of $s+\delta$. Sharpness: $\gamma(s)\le 2\varepsilon$ admits an adversary destroying unique maximality.

\section{Definitions}
\begin{definition}[Unique maximizer]
$i^\star$ uniquely maximizes $s$ if $s_j\le s_{i^\star}$ for all $j$ and $s_j=s_{i^\star}\Rightarrow j=i^\star$.
\end{definition}
\begin{definition}[Margin]
$\gamma(s)=s_{i^\star}-\max_{j\neq i^\star}s_j$.\end{definition}

\section{Theorem}
\begin{theorem}[Margin invariance]\label{thm:inv}
If $i^\star$ uniquely maximizes $s$ and $s_{i^\star}-s_j>2\varepsilon$ for all $j\neq i^\star$, then for every $\|\delta\|_\infty\le\varepsilon$, $i^\star$ uniquely maximizes $s+\delta$.
\end{theorem}
\begin{theorem}[Sharpness]\label{thm:sharp}
If some $j\neq i^\star$ has $s_{i^\star}-s_j\le 2\varepsilon$, an admissible $\delta$ destroys unique maximality.
\end{theorem}

\section{Intuition}
Depressing the winner by $\varepsilon$ and elevating a rival by $\varepsilon$ consumes $2\varepsilon$ of margin. Strict margin above $2\varepsilon$ leaves the winner unique; otherwise a two-point adversary forces a tie or loss.

\section{Examples}
\begin{example}
Scores $(7,5,4)$, $i^\star=0$, $\gamma=2$. For $\varepsilon=0.8$, $2\varepsilon=1.6<2$, so the winner is stable. For $\varepsilon=1$, $2\varepsilon=2\ge\gamma$; $\delta=(-1,+1,0)$ yields a tie between the first two scores.
\end{example}

\section{Proof}
Let $s'=s+\delta$ with $\|\delta\|_\infty\le\varepsilon$. For $j\neq i^\star$,
\[
s'_{i^\star}-s'_j
=(s_{i^\star}-s_j)+(\delta_{i^\star}-\delta_j)
\ge (s_{i^\star}-s_j)-2\varepsilon
>0,
\]
using $s_{i^\star}-s_j>2\varepsilon$. Thus $s'_j<s'_{i^\star}$ for all $j\neq i^\star$, so $i^\star$ uniquely maximizes $s'$ (Theorem~\ref{thm:inv}).

For sharpness, pick $j$ with $s_{i^\star}-s_j\le 2\varepsilon$, set $\delta_{i^\star}=-\varepsilon$, $\delta_j=+\varepsilon$, and $0$ elsewhere. Then $\|\delta\|_\infty\le\varepsilon$ and $s'_{i^\star}-s'_j=\gamma_{ij}-2\varepsilon\le 0$, so unique maximality fails (Theorem~\ref{thm:sharp}).

\section{Formal statement}
{\raggedright\sloppy
\begin{description}[style=nextline,leftmargin=1.1em,font=\normalfont\bfseries]
\item[Lean module] \leanpath{Research.Operators.WeightedTournamentWinner.Margin}
\item[Source file] \leanpath{lean/Research/Operators/WeightedTournamentWinner/Margin.lean}
\item[Theorems]\leavevmode\par\vspace{0.15em}\begin{itemize}[leftmargin=1.2em]
  \item \leanpath{weighted\_tournament\_winner\_margin\_invariance}
  \item \leanpath{weighted\_tournament\_winner\_margin\_sharpness}
\end{itemize}
\item[Note] Definitional reduction of Argmax.Margin.
\item[Certificate] \leanpath{lean/certificates/weighted-tournament-winner/weighted-tournament-winner-margin}
\item[Status] \texttt{LEAN\_FULL} on Mathlib $\mathbb{R}$ (\texttt{REAL\_MATHLIB})
\end{description}
\par}

\section{Proof dependencies}
{\raggedright\sloppy
\begin{itemize}
\item \leanpath{Research.Operators.Argmax.Margin} (\leanpath{margin\_invariance}, \leanpath{margin\_sharpness}).
\item $\ell_\infty$ gap shrinkage by at most $2\varepsilon$.
\item No other operator theorems required.
\end{itemize}
\par}

\section{Consequences}
{\raggedright\sloppy
Any unique-max presentation of \leanpath{weighted-tournament-winner} inherits Argmax margin stability. Related: \leanpath{argmax}, \leanpath{tournament-winner}, \leanpath{greedy-maximum-selection}.
\par}

\end{document}
